% ==============================================================================
% Capítulo de Resultados
% Tabelas geradas por ../Modelagem/regressao/06_tabelas_latex.R e copiadas para
% tabelas/resultados/ (não editar as tabelas à mão).
% ==============================================================================

\chapter{Resultados}\label{cap:resultados}

Este capítulo apresenta os resultados dos dois modelos descritos na \autoref{sec:especificacao}, com a despesa em logaritmo. Os efeitos das dummies são apresentados como variação percentual exata. A leitura segue o critério definido na \autoref{sec:especificacao}: um efeito é considerado sólido quando é significativo tanto com o erro-padrão de Driscoll e Kraay quanto com o erro agrupado por município, quando não depende de uma só eleição e quando vale nas duas formas funcionais.

% ------------------------------------------------------------------------------
\section{Ciclo eleitoral}\label{sec:res_ciclo}

A \autoref{tab:regressao_principal_A} apresenta as estimativas do modelo A. Os resultados indicam que, mantidas constantes as demais variáveis explicativas, a despesa total per capita é aproximadamente 5,3\% superior nos anos eleitorais. Esse efeito agregado, entretanto, encobre uma heterogeneidade considerável entre as funções de gasto. Os maiores efeitos são observados em Habitação e Urbanismo (27\%) e Transporte (19,5\%), seguidos por Assistência Social (7\%) e Saúde e Saneamento (4,2\%), com coeficientes estatisticamente significativos ao nível de 1\%. Para Educação e Cultura, a estimativa corresponde a um aumento de aproximadamente 3,5\%, significativo ao nível de 5\%. Esse resultado, contudo, apresenta menor robustez, uma vez que sua significância estatística depende do tratamento conferido ao ano de 2020 (\autoref{sec:res_robustez}). Em Agricultura, o efeito estimado é de 5,4\%, significativo ao nível de 1\% na especificação logarítmica, mas não estatisticamente significativo na especificação em nível (R3), o que evidencia sua sensibilidade à forma funcional adotada. Para Administração, não se identifica efeito estatisticamente significativo. O coeficiente associado à função Comunicações não é interpretado. Por fim, a magnitude da estimativa para Transporte deve ser considerada à luz da ressalva metodológica apresentada na \autoref{sec:especificacao}.\footnote{Em Transporte, cerca de 13\% dos valores per capita são de até R\$~10, faixa em que o coeficiente em logaritmo não se lê como variação percentual.}

Esse padrão é consistente com a hipótese de que os efeitos do ciclo eleitoral se manifestam predominantemente na composição funcional da despesa pública, por meio da realocação de recursos em favor de áreas potencialmente mais visíveis ao eleitorado, como obras urbanas e transporte \parencite{drazen2010}. Essa interpretação também encontra respaldo no comportamento apresentado na \autoref{fig:grupos1}, que evidencia um padrão de dente de serra nas despesas com Habitação e Urbanismo e Transporte, caracterizado por elevações nos anos eleitorais e retrações no período subsequente.

No ano pré-eleitoral, por sua vez, não se observa efeito estatisticamente significativo sobre a despesa total per capita. Identificam-se, contudo, efeitos positivos em algumas funções específicas: Educação e Cultura (4,7\%), Habitação e Urbanismo (5,2\%) e Agricultura (5,8\%), com significância estatística ao nível de 5\%, e Assistência Social (3,6\%), significativa ao nível de 10\%. Esses resultados sugerem que a antecipação do ciclo eleitoral pode estar associada a alterações seletivas na alocação das despesas, embora sem produzir uma expansão estatisticamente identificável do gasto agregado.

Por seu turno, a variável indicadora da pandemia de covid-19 permite identificar as particularidades de 2020 em relação aos demais anos da amostra. As estimativas apontam para uma despesa aproximadamente 9,5\% superior em Saúde e Saneamento e 14,9\% inferior em Educação e Cultura, resultados compatíveis, respectivamente, com a intensificação da demanda por serviços de saúde e a suspensão das atividades escolares presenciais naquele ano. Para a despesa total per capita, entretanto, o coeficiente associado à pandemia não apresenta significância estatística.

Entre as variáveis de controle, os coeficientes associados às receitas podem ser interpretados como elasticidades. Para a despesa total per capita, as estimativas indicam que um aumento de 1\% na receita tributária está associado a um acréscimo de aproximadamente 0,084\% na despesa, enquanto um aumento de 1\% nas transferências correntes está associado a uma elevação de aproximadamente 0,32\%. O coeficiente do PIB nacional, por sua vez, não apresenta significância estatística. Conforme discutido na \autoref{sec:variaveis}, os coeficientes das variáveis provenientes do Censo não são objeto de interpretação.

No que se refere à literatura, os resultados corroboram parcialmente os achados de \textcite{sakurai2009}, que analisa o período de 1990 a 2005 e também identifica efeitos positivos do ano eleitoral sobre as despesas com Saúde e Saneamento, Habitação e Urbanismo e Transporte. Embora o autor utilize uma classificação que agrega Assistência e Previdência, a estimação com esse mesmo agrupamento também revela, neste trabalho, um coeficiente positivo e estatisticamente significativo ao nível de 1\% para o ano eleitoral (\autoref{tab:robustez_R8}, no \autoref{ap:robustez}). As divergências concentram-se em duas funções. Em Agricultura, \textcite{sakurai2009} encontra uma redução da despesa nos anos eleitorais, enquanto, neste estudo, o coeficiente é positivo na especificação logarítmica, mas não apresenta significância estatística na especificação em nível. Em Educação e Cultura, o autor não identifica efeito eleitoral estatisticamente significativo, ao passo que os resultados aqui obtidos apontam para um efeito positivo, embora de menor magnitude e robustez.

% Gerado por regressao/06_tabelas_latex.R (projeto Modelagem). Não editar à mão.
\begin{sidewaystable}
\caption{Resultados do modelo A (ciclo eleitoral) – efeitos fixos de município}\label{tab:regressao_principal_A}
\centering
\footnotesize
\setlength{\tabcolsep}{3pt}
\renewcommand{\arraystretch}{0.9}
\begin{tabular}{@{}>{\raggedright\arraybackslash}p{\dimexpr\linewidth-594pt\relax}>{\centering\arraybackslash}p{60pt}>{\centering\arraybackslash}p{60pt}>{\centering\arraybackslash}p{60pt}>{\centering\arraybackslash}p{60pt}>{\centering\arraybackslash}p{60pt}>{\centering\arraybackslash}p{60pt}>{\centering\arraybackslash}p{60pt}>{\centering\arraybackslash}p{60pt}>{\centering\arraybackslash}p{60pt}@{}}
\hline
\textbf{Variável} & \textbf{\hspace{0pt}Despesa total} & \textbf{\hspace{0pt}Saúde e Saneamento} & \textbf{\hspace{0pt}Educação e Cultura} & \textbf{\hspace{0pt}Habitação e Urbanismo} & \textbf{\hspace{0pt}Assistência Social} & \textbf{\hspace{0pt}Transporte} & \textbf{\hspace{0pt}Administração} & \textbf{\hspace{0pt}Agricultura} & \textbf{\hspace{0pt}Comunicações} \\ \hline
Ano eleitoral & 0,0519\textsuperscript{***} & 0,0411\textsuperscript{***} & 0,0343\textsuperscript{**} & 0,2389\textsuperscript{***} & 0,0678\textsuperscript{***} & 0,1778\textsuperscript{***} & 0,0097 & 0,0529\textsuperscript{***} & 0,0318\textsuperscript{***} \\
 & (0,0088) & (0,0044) & (0,0149) & (0,0161) & (0,0113) & (0,0198) & (0,0077) & (0,0142) & (0,0106) \\
Ano pré-eleitoral & 0,0209 & $-$0,0083 & 0,0461\textsuperscript{**} & 0,0509\textsuperscript{**} & 0,0355\textsuperscript{*} & 0,0188 & 0,0077 & 0,0564\textsuperscript{**} & 0,0605\textsuperscript{***} \\
 & (0,0143) & (0,0067) & (0,0228) & (0,0236) & (0,0198) & (0,0239) & (0,0086) & (0,0224) & (0,0122) \\
Pandemia (2020) & $-$0,0058 & 0,0908\textsuperscript{***} & $-$0,1610\textsuperscript{***} & 0,1254\textsuperscript{***} & 0,0169\textsuperscript{*} & 0,0711\textsuperscript{***} & 0,0297\textsuperscript{***} & 0,0312\textsuperscript{**} & 0,0444\textsuperscript{***} \\
 & (0,0072) & (0,0035) & (0,0125) & (0,0117) & (0,0094) & (0,0196) & (0,0048) & (0,0142) & (0,0130) \\
Prefeito na coalizão & 0,0035\textsuperscript{**} & 0,0012 & 0,0032 & 0,0155\textsuperscript{***} & 0,0010 & 0,0111 & 0,0066\textsuperscript{***} & 0,0113\textsuperscript{*} & $-$0,0019 \\
do governador & (0,0018) & (0,0023) & (0,0028) & (0,0051) & (0,0035) & (0,0095) & (0,0018) & (0,0067) & (0,0085) \\
Prefeito na coalizão & 0,0092\textsuperscript{*} & 0,0058 & 0,0142\textsuperscript{*} & $-$0,0039 & 0,0111 & 0,0461\textsuperscript{***} & 0,0076\textsuperscript{**} & 0,0172\textsuperscript{**} & 0,0155 \\
do presidente & (0,0050) & (0,0036) & (0,0082) & (0,0082) & (0,0078) & (0,0113) & (0,0033) & (0,0081) & (0,0230) \\
Receita tributária & 0,0837\textsuperscript{***} & 0,0605\textsuperscript{***} & 0,0840\textsuperscript{***} & 0,1639\textsuperscript{***} & 0,0560\textsuperscript{***} & 0,1606\textsuperscript{***} & 0,0730\textsuperscript{***} & 0,0865\textsuperscript{***} & 0,0830\textsuperscript{***} \\
\textit{per capita} (log) & (0,0086) & (0,0063) & (0,0109) & (0,0118) & (0,0048) & (0,0231) & (0,0075) & (0,0117) & (0,0214) \\
Transferências correntes & 0,3203\textsuperscript{***} & 0,2908\textsuperscript{***} & 0,3476\textsuperscript{***} & 0,4887\textsuperscript{***} & 0,2860\textsuperscript{***} & 0,5435\textsuperscript{***} & 0,2626\textsuperscript{***} & 0,3724\textsuperscript{***} & 0,3244\textsuperscript{***} \\
\textit{per capita} (log) & (0,0577) & (0,0532) & (0,0702) & (0,0854) & (0,0479) & (0,0840) & (0,0341) & (0,0601) & (0,0835) \\
Proporção de jovens & $-$0,6239 & $-$1,5188\textsuperscript{***} & 0,4320 & $-$4,1290\textsuperscript{***} & 0,0340 & $-$6,2984\textsuperscript{***} & $-$1,2191\textsuperscript{***} & $-$4,4731\textsuperscript{***} & $-$3,4586\textsuperscript{***} \\
 & (0,6429) & (0,5026) & (0,9373) & (1,3463) & (0,6465) & (2,2367) & (0,4658) & (0,9064) & (0,7112) \\
Proporção de idosos & $-$2,2218\textsuperscript{***} & $-$3,3555\textsuperscript{***} & $-$3,0263\textsuperscript{***} & 2,8280\textsuperscript{**} & $-$1,9011\textsuperscript{**} & 2,6502 & $-$2,2522\textsuperscript{***} & $-$1,0664 & $-$8,2181\textsuperscript{***} \\
 & (0,6417) & (0,6651) & (1,0569) & (1,4127) & (0,9505) & (1,7568) & (0,3411) & (1,9546) & (1,6829) \\
Grau de urbanização & 0,1225\textsuperscript{**} & 0,0360 & 0,0451 & 0,5168\textsuperscript{***} & 0,0665 & 0,0499 & 0,0553\textsuperscript{**} & 0,5959\textsuperscript{***} & $-$0,5338 \\
 & (0,0487) & (0,0673) & (0,0579) & (0,1088) & (0,0717) & (0,1864) & (0,0252) & (0,1769) & (0,3342) \\
População (log) & $-$0,5413\textsuperscript{***} & $-$0,6055\textsuperscript{***} & $-$0,5020\textsuperscript{***} & $-$0,4093\textsuperscript{***} & $-$0,6695\textsuperscript{***} & $-$0,2293\textsuperscript{*} & $-$0,5539\textsuperscript{***} & $-$0,4928\textsuperscript{***} & $-$0,5490\textsuperscript{***} \\
 & (0,0852) & (0,0622) & (0,1162) & (0,1455) & (0,0728) & (0,1191) & (0,0610) & (0,0953) & (0,1364) \\
PIB nacional real (log) & 0,6812 & 1,2351\textsuperscript{***} & $-$0,4691 & 3,2346\textsuperscript{***} & 1,0233\textsuperscript{*} & 2,5929\textsuperscript{***} & 0,5371\textsuperscript{*} & 2,0275\textsuperscript{***} & 1,0143\textsuperscript{***} \\
 & (0,4314) & (0,1457) & (0,7428) & (0,6231) & (0,5383) & (0,6719) & (0,2742) & (0,7220) & (0,3843) \\
Tendência & 0,0009 & 0,0361\textsuperscript{***} & $-$0,0413 & 0,0304 & 0,0253 & $-$0,0482 & $-$0,0163 & $-$0,0229 & $-$0,0306 \\
 & (0,0206) & (0,0088) & (0,0344) & (0,0295) & (0,0261) & (0,0399) & (0,0134) & (0,0401) & (0,0228) \\
Tendência ao quadrado & 0,0014 & $-$0,0014\textsuperscript{*} & 0,0060\textsuperscript{*} & $-$0,0062\textsuperscript{**} & $-$0,0009 & $-$0,0008 & 0,0014 & $-$0,0006 & $-$0,0005 \\
 & (0,0020) & (0,0008) & (0,0032) & (0,0029) & (0,0024) & (0,0033) & (0,0013) & (0,0034) & (0,0017) \\
\hline
Observações & 70.669 & 70.526 & 70.550 & 69.574 & 70.484 & 54.422 & 70.553 & 64.464 & 12.360 \\
Municípios & 5.568 & 5.568 & 5.568 & 5.564 & 5.568 & 5.205 & 5.568 & 5.436 & 1.896 \\
R\textsuperscript{2} & 0,851 & 0,680 & 0,680 & 0,264 & 0,379 & 0,137 & 0,306 & 0,134 & 0,031 \\
\hline
\end{tabular}
\fonte{Elaboração própria.}
\nota{Variáveis dependentes: despesa por grupo de funções. Despesa e receitas \textit{per capita}, em R\$ de 2025, em logaritmo (\autoref{sec:especificacao}). O coeficiente de uma \textit{dummy} se lê como variação percentual aproximada (exata: 100 $\times$ (exp(b) $-$ 1)). PIB nacional em log. Efeito fixo de município. Erros-padrão de Driscoll-Kraay entre parênteses. \textsuperscript{***}p$<$0,01; \textsuperscript{**}p$<$0,05; \textsuperscript{*}p$<$0,1. Os coeficientes de coalizão são apresentados para completude; a análise das coalizões usa o modelo B.}
\end{sidewaystable}


% ------------------------------------------------------------------------------
\section{Efeito por eleição}\label{sec:res_por_eleicao}

O coeficiente do ano eleitoral estimado no modelo A corresponde a uma média dos efeitos das eleições de 2016 e 2024, uma vez que o ano de 2020 é absorvido pela variável indicadora da pandemia. Para verificar se esse efeito médio é homogêneo entre os pleitos, a \autoref{tab:robustez_R16} estima um coeficiente para cada eleição e para cada ano pré-eleitoral, e os resultados revelam diferenças entre as duas eleições.

Para a despesa total per capita, mantidas constantes as demais variáveis, o efeito estimado é de aproximadamente 3,3\% em 2016, significativo ao nível de 5\%, e de 6,7\% em 2024, significativo ao nível de 1\%. O efeito de 2016, além de menor, é sensível à especificação: no teste R7a, em que a amostra se encerra em 2023 e apenas a eleição de 2016 identifica o efeito, o coeficiente da despesa total não apresenta significância estatística (\autoref{tab:robustez_R7a}, no \autoref{ap:robustez}); quando, além do corte da amostra, o PIB nacional é substituído pelo PIB municipal (R7), o coeficiente é de aproximadamente 2,7\%, significativo ao nível de 1\% (\autoref{tab:robustez_R7}).

Nas funções de gasto, o padrão de composição identificado na \autoref{sec:res_ciclo} repete-se nos dois pleitos. Habitação e Urbanismo (aproximadamente 24\% em 2016 e 28,7\% em 2024), Transporte (21,2\% e 17,6\%) e Saúde e Saneamento (4,2\% e 4\%) apresentam aumentos nas duas eleições, todos significativos ao nível de 1\%, o que indica que esses efeitos não dependem de um único pleito. Em Administração, o coeficiente de 2016 não é estatisticamente significativo, e Assistência Social e Educação e Cultura apresentam efeito significativo apenas em 2024.

Nos anos pré-eleitorais de 2015 e 2019, não se observa efeito estatisticamente significativo sobre a despesa total, embora algumas funções apresentem variações: Saúde e Saneamento registra reduções nos dois anos (aproximadamente 1,1\% em 2015 e 1,9\% em 2019), e Habitação e Urbanismo, aumento em 2015 (4,3\%), todos significativos ao nível de 1\%.

O ano de 2024 distingue-se dos demais: quase todas as funções apresentam aumento simultâneo, inclusive Administração (aproximadamente 2,2\%, significativo ao nível de 5\%). De modo semelhante, no ano pré-eleitoral de 2023, a despesa total é cerca de 5,2\% superior, com significância ao nível de 5\%, acompanhada de aumentos em várias funções. Esse resultado admite duas leituras. De um lado, uma expansão generalizada, que alcança inclusive o gasto menos visível ao eleitor, é compatível com o forte crescimento das receitas municipais no período, discutido na \autoref{sec:descritivas}. Como os modelos controlam a receita tributária e as transferências correntes, parte desse movimento é absorvida pelas variáveis de receita, e o restante pode refletir outros fatores próprios do período 2023--2024 não captados pelo modelo. De outro lado, os investimentos municipais cresceram 26,1\% em termos reais em 2024, alta que o Ipea associa ao ano eleitoral \parencite[p.~10]{santana2025} e que o BNDES relaciona às eleições locais \parencite[p.~121]{bndes2025}, o que é compatível com uma estratégia eleitoral. O desenho empírico, contudo, não permite separar essas duas explicações, uma vez que tais fatores atingem todos os municípios no mesmo ano.

A hipótese de recomposição do gasto é avaliada diretamente pelo teste R4 (\autoref{tab:robustez_R4}, no \autoref{ap:robustez}), no qual a variável dependente é a participação de cada grupo na despesa total. No ano eleitoral, mantidas constantes as demais variáveis, Habitação e Urbanismo ganha aproximadamente 1,46 ponto percentual de participação, e Transporte, 0,44 ponto, ao passo que Administração perde cerca de 0,45 ponto, e Educação e Cultura, 0,55 ponto, todos com significância ao nível de 1\%. Esse padrão é compatível com o ciclo de composição descrito por \textcite{drazen2010} e com os resultados de \textcite{teixeira2021}, que medem a participação das despesas administrativas sob maior competição eleitoral. Assim, a evidência mais sólida de ciclo eleitoral parece residir na composição do gasto, presente nas duas eleições, enquanto o efeito sobre a despesa total é mais pronunciado em 2024.

% Gerado por regressao/06_tabelas_latex.R (projeto Modelagem). Não editar à mão.
\begin{sidewaystable}
\caption{Robustez R16 — ano eleitoral e pré-eleitoral separados por eleição (modelo A)}\label{tab:robustez_R16}
\centering
\footnotesize
\setlength{\tabcolsep}{3pt}
\begin{tabular}{@{}>{\raggedright\arraybackslash}p{\dimexpr\linewidth-594pt\relax}>{\centering\arraybackslash}p{60pt}>{\centering\arraybackslash}p{60pt}>{\centering\arraybackslash}p{60pt}>{\centering\arraybackslash}p{60pt}>{\centering\arraybackslash}p{60pt}>{\centering\arraybackslash}p{60pt}>{\centering\arraybackslash}p{60pt}>{\centering\arraybackslash}p{60pt}>{\centering\arraybackslash}p{60pt}@{}}
\hline
\textbf{Variável} & \textbf{\hspace{0pt}Despesa total} & \textbf{\hspace{0pt}Saúde e Saneamento} & \textbf{\hspace{0pt}Educação e Cultura} & \textbf{\hspace{0pt}Habitação e Urbanismo} & \textbf{\hspace{0pt}Assistência Social} & \textbf{\hspace{0pt}Transporte} & \textbf{\hspace{0pt}Administração} & \textbf{\hspace{0pt}Agricultura} & \textbf{\hspace{0pt}Comunicações} \\ \hline
\multicolumn{10}{@{}l}{\textit{Modelo A (ciclo eleitoral)}} \\
Ano eleitoral de 2016 & 0,0329\textsuperscript{**} & 0,0414\textsuperscript{***} & 0,0030 & 0,2153\textsuperscript{***} & 0,0402 & 0,1920\textsuperscript{***} & $-$0,0059 & 0,0632\textsuperscript{**} & 0,0393\textsuperscript{***} \\
 & (0,0148) & (0,0086) & (0,0261) & (0,0206) & (0,0256) & (0,0221) & (0,0086) & (0,0256) & (0,0126) \\
Ano eleitoral de 2024 & 0,0650\textsuperscript{***} & 0,0396\textsuperscript{***} & 0,0564\textsuperscript{**} & 0,2526\textsuperscript{***} & 0,0873\textsuperscript{***} & 0,1625\textsuperscript{***} & 0,0217\textsuperscript{**} & 0,0418 & 0,0250 \\
 & (0,0159) & (0,0049) & (0,0246) & (0,0326) & (0,0171) & (0,0396) & (0,0108) & (0,0319) & (0,0163) \\
Ano pré-eleitoral & 0,0083 & $-$0,0106\textsuperscript{***} & 0,0243\textsuperscript{**} & 0,0422\textsuperscript{***} & 0,0221\textsuperscript{*} & $-$0,0070 & $-$0,0066\textsuperscript{**} & 0,0657\textsuperscript{***} & 0,0627\textsuperscript{***} \\
de 2015 & (0,0061) & (0,0039) & (0,0121) & (0,0042) & (0,0120) & (0,0052) & (0,0029) & (0,0097) & (0,0080) \\
Ano pré-eleitoral & $-$0,0025 & $-$0,0196\textsuperscript{***} & 0,0137 & $-$0,0030 & 0,0063 & $-$0,0093 & 0,0024 & 0,0265 & 0,0627\textsuperscript{***} \\
de 2019 & (0,0124) & (0,0061) & (0,0204) & (0,0205) & (0,0197) & (0,0239) & (0,0072) & (0,0223) & (0,0181) \\
Ano pré-eleitoral & 0,0511\textsuperscript{**} & 0,0033 & 0,0920\textsuperscript{**} & 0,1012\textsuperscript{**} & 0,0705\textsuperscript{**} & 0,0695 & 0,0251\textsuperscript{*} & 0,0720 & 0,0566\textsuperscript{***} \\
de 2023 & (0,0238) & (0,0080) & (0,0375) & (0,0443) & (0,0277) & (0,0523) & (0,0152) & (0,0455) & (0,0179) \\
\hline
Observações & 70.669 & 70.526 & 70.550 & 69.574 & 70.484 & 54.422 & 70.553 & 64.464 & 12.360 \\
Municípios & 5.568 & 5.568 & 5.568 & 5.564 & 5.568 & 5.205 & 5.568 & 5.436 & 1.896 \\
R\textsuperscript{2} & 0,853 & 0,680 & 0,683 & 0,265 & 0,380 & 0,137 & 0,306 & 0,134 & 0,031 \\
\hline
\end{tabular}
\fonte{Elaboração própria.}
\nota{Despesa e receitas \textit{per capita}, em R\$ de 2025, em logaritmo (\autoref{sec:especificacao}). O coeficiente de uma \textit{dummy} se lê como variação percentual aproximada (exata: 100 $\times$ (exp(b) $-$ 1)). Só modelo A: efeito fixo de município; ano eleitoral e pré-eleitoral separados por eleição (2020 absorvido pela \textit{dummy} de pandemia). Erros-padrão de Driscoll-Kraay entre parênteses. \textsuperscript{***}p$<$0,01; \textsuperscript{**}p$<$0,05; \textsuperscript{*}p$<$0,1.}
\end{sidewaystable}


% ------------------------------------------------------------------------------
\section{Coalizões}\label{sec:res_coalizoes}

A \autoref{tab:regressao_principal_B} apresenta as estimativas do modelo B, que compara prefeitos alinhados e não alinhados no mesmo ano. Os resultados indicam que, mantidas constantes as demais variáveis, os municípios cujo prefeito integra a coligação do governador apresentam despesa total per capita aproximadamente 0,37\% superior, com significância ao nível de 1\%. Trata-se de um efeito de pequena magnitude, mas estável: o coeficiente permanece significativo ao nível de 1\% em todos os testes do modelo B que incluem a despesa total, com valores entre 0,0030 e 0,0038, e também na especificação em nível, na qual corresponde a R\$~31,56 per capita (\autoref{tab:robustez_R3}). Em Administração, o efeito é de aproximadamente 0,7\%, significativo ao nível de 1\% nas duas formas funcionais (R\$~7,17 em nível). Em Habitação e Urbanismo, o efeito é de cerca de 1,6\%, significativo ao nível de 5\%, porém apenas na especificação logarítmica. Esses resultados são compatíveis com os de \textcite{ferreira2007}, segundo os quais municípios governados por aliados do governador recebem mais transferências voluntárias. Como as transferências correntes já figuram entre os controles, o efeito estimado pode indicar um canal adicional, possivelmente associado a transferências de capital ou a outras formas de apoio estadual não captadas pelas transferências correntes.

O alinhamento com o presidente, por outro lado, não apresenta efeito estatisticamente significativo sobre a despesa total na amostra completa. Na amostra de 2013 a 2023, entretanto, o coeficiente é negativo e significativo ao nível de 5\%, de aproximadamente $-$0,3\% (R7 e R7a; \autoref{tab:robustez_R7} e \autoref{tab:robustez_R7a}), o que indica que o resultado é sensível ao período considerado. Em Transporte, o coeficiente é positivo na especificação logarítmica e negativo na especificação em nível, ambos significativos ao nível de 1\%; por depender da forma funcional, esse resultado não é interpretado. Em outras funções, a significância dos coeficientes também difere entre as duas formas funcionais (\autoref{tab:regressao_principal_B} e \autoref{tab:robustez_R3}), razão pela qual esses resultados não são destacados.

A comparação com o modelo A evidencia a necessidade do efeito fixo de ano. Na ausência desse controle, a coalizão com o presidente aparece associada a uma despesa total aproximadamente 0,9\% superior, significativa ao nível de 10\% (\autoref{tab:regressao_principal_A}); com ele, o coeficiente aproxima-se de zero e não apresenta significância estatística. Como a proporção de prefeitos alinhados ao presidente varia de forma abrupta entre os governos federais (\autoref{tab:politicas}), o coeficiente do modelo A tende a confundir o efeito do alinhamento com as diferenças entre os períodos. Já o efeito do alinhamento com o governador, cuja variação ocorre entre estados no mesmo ano, mantém-se nos dois modelos e também com o erro-padrão de Driscoll e Kraay (\autoref{tab:robustez_R13}, no \autoref{ap:robustez}).

No que se refere à literatura, \textcite{sakurai2009} identifica efeitos do alinhamento com o governador em diversas funções, com sinais distintos: positivos em Saúde e Saneamento e em Habitação e Urbanismo, e negativos em Educação e Cultura, Agricultura e Assistência e Previdência. Neste trabalho, com efeito fixo de ano, o efeito é positivo na despesa total e em Administração, nas duas formas funcionais.

% Gerado por regressao/06_tabelas_latex.R (projeto Modelagem). Não editar à mão.
\begin{sidewaystable}
\caption{Resultados do modelo B (coalizões) – efeitos fixos de município e de ano}\label{tab:regressao_principal_B}
\centering
\footnotesize
\setlength{\tabcolsep}{3pt}
\begin{tabular}{@{}>{\raggedright\arraybackslash}p{\dimexpr\linewidth-594pt\relax}>{\centering\arraybackslash}p{60pt}>{\centering\arraybackslash}p{60pt}>{\centering\arraybackslash}p{60pt}>{\centering\arraybackslash}p{60pt}>{\centering\arraybackslash}p{60pt}>{\centering\arraybackslash}p{60pt}>{\centering\arraybackslash}p{60pt}>{\centering\arraybackslash}p{60pt}>{\centering\arraybackslash}p{60pt}@{}}
\hline
\textbf{Variável} & \textbf{\hspace{0pt}Despesa total} & \textbf{\hspace{0pt}Saúde e Saneamento} & \textbf{\hspace{0pt}Educação e Cultura} & \textbf{\hspace{0pt}Habitação e Urbanismo} & \textbf{\hspace{0pt}Assistência Social} & \textbf{\hspace{0pt}Transporte} & \textbf{\hspace{0pt}Administração} & \textbf{\hspace{0pt}Agricultura} & \textbf{\hspace{0pt}Comunicações} \\ \hline
Prefeito na coalizão & 0,0037\textsuperscript{***} & 0,0009 & 0,0033\textsuperscript{**} & 0,0160\textsuperscript{**} & 0,0009 & 0,0114 & 0,0070\textsuperscript{***} & 0,0110 & $-$0,0022 \\
do governador & (0,0010) & (0,0018) & (0,0017) & (0,0064) & (0,0029) & (0,0100) & (0,0027) & (0,0075) & (0,0175) \\
Prefeito na coalizão & $-$0,0010 & 0,0001 & $-$0,0025 & $-$0,0163\textsuperscript{*} & $-$0,0058 & 0,0418\textsuperscript{***} & 0,0043 & 0,0045 & 0,0075 \\
do presidente & (0,0014) & (0,0025) & (0,0021) & (0,0085) & (0,0040) & (0,0138) & (0,0037) & (0,0102) & (0,0243) \\
Receita tributária & 0,0776\textsuperscript{***} & 0,0583\textsuperscript{***} & 0,0742\textsuperscript{***} & 0,1550\textsuperscript{***} & 0,0471\textsuperscript{***} & 0,1521\textsuperscript{***} & 0,0692\textsuperscript{***} & 0,0775\textsuperscript{***} & 0,0786\textsuperscript{***} \\
per capita (log) & (0,0037) & (0,0043) & (0,0048) & (0,0123) & (0,0059) & (0,0188) & (0,0057) & (0,0140) & (0,0289) \\
Transferências correntes & 0,3164\textsuperscript{***} & 0,2905\textsuperscript{***} & 0,3427\textsuperscript{***} & 0,4729\textsuperscript{***} & 0,2860\textsuperscript{***} & 0,5180\textsuperscript{***} & 0,2595\textsuperscript{***} & 0,3549\textsuperscript{***} & 0,3247\textsuperscript{***} \\
per capita (log) & (0,0351) & (0,0340) & (0,0380) & (0,0577) & (0,0326) & (0,0849) & (0,0322) & (0,0452) & (0,1065) \\
Proporção de jovens & $-$0,5308\textsuperscript{***} & $-$1,5562\textsuperscript{***} & 0,4248\textsuperscript{*} & $-$3,2494\textsuperscript{***} & $-$0,1301 & $-$5,0797\textsuperscript{***} & $-$1,1059\textsuperscript{***} & $-$4,0292\textsuperscript{***} & $-$3,7329\textsuperscript{*} \\
 & (0,1195) & (0,2015) & (0,2324) & (0,7680) & (0,3308) & (1,1141) & (0,3123) & (0,8256) & (2,2294) \\
Proporção de idosos & $-$2,3573\textsuperscript{***} & $-$3,3172\textsuperscript{***} & $-$3,0680\textsuperscript{***} & 1,7387 & $-$1,7420\textsuperscript{***} & 1,1688 & $-$2,4043\textsuperscript{***} & $-$1,6408 & $-$7,8752\textsuperscript{**} \\
 & (0,2134) & (0,3181) & (0,2760) & (1,2144) & (0,5369) & (1,5377) & (0,4601) & (1,1414) & (3,4568) \\
Grau de urbanização & 0,1178\textsuperscript{***} & 0,0363 & 0,0416 & 0,4843\textsuperscript{**} & 0,0691 & $-$0,0069 & 0,0512 & 0,5688\textsuperscript{**} & $-$0,5211 \\
 & (0,0395) & (0,0565) & (0,0893) & (0,2447) & (0,1125) & (0,3826) & (0,1165) & (0,2754) & (0,8577) \\
População (log) & $-$0,4974\textsuperscript{***} & $-$0,5871\textsuperscript{***} & $-$0,4299\textsuperscript{***} & $-$0,3514\textsuperscript{***} & $-$0,6097\textsuperscript{***} & $-$0,1743 & $-$0,5304\textsuperscript{***} & $-$0,4347\textsuperscript{***} & $-$0,4904\textsuperscript{**} \\
 & (0,0313) & (0,0329) & (0,0358) & (0,0926) & (0,0418) & (0,1357) & (0,0399) & (0,0931) & (0,2229) \\
\hline
Observações & 70.669 & 70.526 & 70.550 & 69.574 & 70.484 & 54.422 & 70.553 & 64.464 & 12.360 \\
Municípios & 5.568 & 5.568 & 5.568 & 5.564 & 5.568 & 5.205 & 5.568 & 5.436 & 1.896 \\
R\textsuperscript{2} & 0,390 & 0,163 & 0,196 & 0,036 & 0,069 & 0,019 & 0,075 & 0,017 & 0,013 \\
\hline
\end{tabular}
\fonte{Elaboração própria.}
\nota{Variáveis dependentes: despesa por grupo de funções. Despesa e receitas per capita, em R\$ de 2025, em logaritmo (\autoref{sec:especificacao}). O coeficiente de uma dummy se lê como variação percentual aproximada (exata: 100 $\times$ (exp(b) $-$ 1)). Efeitos fixos de município e de ano. Erros-padrão agrupados por município entre parênteses. \textsuperscript{***}p$<$0,01; \textsuperscript{**}p$<$0,05; \textsuperscript{*}p$<$0,1.}
\end{sidewaystable}


% ------------------------------------------------------------------------------
\section{Robustez}\label{sec:res_robustez}

A \autoref{tab:robustez_resumo} compara os efeitos do ano eleitoral e do ano pré-eleitoral estimados no modelo principal com os obtidos em quatro especificações alternativas.

Na especificação de \textcite{sakurai2009} (R5), que não inclui a dummy de pandemia e utiliza erro-padrão agrupado por município, os efeitos em Habitação e Urbanismo, Transporte, Assistência Social e Saúde e Saneamento permanecem positivos e estatisticamente significativos. A principal diferença ocorre em Educação e Cultura, cujo coeficiente se torna negativo (aproximadamente $-$2,5\%, significativo ao nível de 1\%). O teste R14, que retira apenas a dummy de pandemia, contribui para compreender essa diferença: sem ela, o coeficiente de Educação e Cultura aproxima-se de zero e deixa de ser significativo, pois a queda desse gasto em 2020 passa a ser atribuída à eleição. O resultado de Educação e Cultura depende, portanto, do tratamento conferido a 2020 e não deve ser interpretado como evidência de ciclo eleitoral.

No que diz respeito ao tipo de erro-padrão, a estimação com erro agrupado por município (R12) preserva os coeficientes do modelo principal, mas produz, em geral, erros-padrão menores, de modo que todos os coeficientes do ano eleitoral e do ano pré-eleitoral se tornam significativos. Esse resultado indica que o erro-padrão de Driscoll e Kraay constitui, na maior parte dos casos, a escolha mais conservadora para variáveis comuns a todos os municípios.

Quanto à forma funcional, na especificação em nível (R3), o ano eleitoral é significativo ao nível de 1\% na despesa total (R\$~277,87 per capita, cerca de 5\% da média da \autoref{tab:desc_dependentes}), em Habitação e Urbanismo (R\$~126,81, cerca de 27\%), em Saúde e Saneamento (R\$~53,16, cerca de 4\%), em Transporte (R\$~41,74, cerca de 16\%) e em Assistência Social (R\$~14,08, cerca de 7\%), e significativo ao nível de 10\% em Educação e Cultura. Convertidas em proporção da média, essas magnitudes são da mesma ordem de grandeza das estimadas em logaritmo. Administração e Agricultura não apresentam significância estatística. No ano pré-eleitoral, a despesa total passa a ser significativa ao nível de 10\%, e Agricultura deixa de ser significativa.

Os demais testes, no \autoref{ap:robustez}, avaliam a padronização das siglas, as eleições suplementares, a participação de cada área na despesa total, o PIB municipal, a limpeza da base, o registro irregular de algumas funções, os controles do Censo e a série de população. De modo geral, esses testes preservam os principais resultados, com algumas exceções. Na amostra de 2013 a 2023, o alinhamento com o presidente fica negativo na despesa total, o ano eleitoral deixa de ser significativo em Assistência Social e em Agricultura (e, no R7a, também em Educação e Cultura), os efeitos do ano pré-eleitoral em Educação e Cultura, Habitação e Urbanismo, Assistência Social e Agricultura também deixam de ser significativos, e o do ano pré-eleitoral em Saúde e Saneamento passa a ser significativo e negativo. Com o PIB municipal (R7), o efeito do ano eleitoral em Saúde e Saneamento deixa de ser significativo, e o de Administração torna-se negativo. Por fim, o alinhamento com o governador em Educação e Cultura deixa de ser significativo no R13.

Em síntese, os resultados indicam que: (i) Habitação e Urbanismo, Transporte e Saúde e Saneamento sobem no ano eleitoral nas duas eleições, nas duas formas funcionais e com os dois erros-padrão; (ii) o alinhamento com o governador aumenta a despesa total, com efeito pequeno, de cerca de 0,4\%; (iii) o alinhamento com o presidente não tem efeito sobre a despesa total na amostra completa, e na amostra até 2023 o efeito é negativo, de modo que o resultado é sensível ao período; (iv) a despesa total sobe cerca de 5\% no ano eleitoral, mais em 2024, e em 2016 o efeito é menor e não resiste ao corte da amostra na especificação principal (R7a); e (v) alguns resultados dependem da forma funcional, como Agricultura no ano eleitoral e Transporte e Habitação e Urbanismo nas coalizões, e Comunicações não é interpretada.

% Gerado por regressao/06_tabelas_latex.R (projeto Modelagem). Não editar à mão.
\begin{sidewaystable}
\caption{Robustez do ciclo eleitoral: modelo principal e testes selecionados}\label{tab:robustez_resumo}
\centering
\footnotesize
\setlength{\tabcolsep}{3pt}
\begin{tabular}{@{}>{\raggedright\arraybackslash}p{\dimexpr\linewidth-594pt\relax}>{\centering\arraybackslash}p{60pt}>{\centering\arraybackslash}p{60pt}>{\centering\arraybackslash}p{60pt}>{\centering\arraybackslash}p{60pt}>{\centering\arraybackslash}p{60pt}>{\centering\arraybackslash}p{60pt}>{\centering\arraybackslash}p{60pt}>{\centering\arraybackslash}p{60pt}>{\centering\arraybackslash}p{60pt}@{}}
\hline
\textbf{Variável} & \textbf{\hspace{0pt}Despesa total} & \textbf{\hspace{0pt}Saúde e Saneamento} & \textbf{\hspace{0pt}Educação e Cultura} & \textbf{\hspace{0pt}Habitação e Urbanismo} & \textbf{\hspace{0pt}Assistência Social} & \textbf{\hspace{0pt}Transporte} & \textbf{\hspace{0pt}Administração} & \textbf{\hspace{0pt}Agricultura} & \textbf{\hspace{0pt}Comunicações} \\ \hline
\multicolumn{10}{@{}l}{\textit{Modelo principal (erro de Driscoll-Kraay)}} \\
Ano eleitoral & 0,0519\textsuperscript{***} & 0,0411\textsuperscript{***} & 0,0343\textsuperscript{**} & 0,2389\textsuperscript{***} & 0,0678\textsuperscript{***} & 0,1778\textsuperscript{***} & 0,0097 & 0,0529\textsuperscript{***} & 0,0318\textsuperscript{***} \\
 & (0,0088) & (0,0044) & (0,0149) & (0,0161) & (0,0113) & (0,0198) & (0,0077) & (0,0142) & (0,0106) \\
Ano pré-eleitoral & 0,0209 & $-$0,0083 & 0,0461\textsuperscript{**} & 0,0509\textsuperscript{**} & 0,0355\textsuperscript{*} & 0,0188 & 0,0077 & 0,0564\textsuperscript{**} & 0,0605\textsuperscript{***} \\
 & (0,0143) & (0,0067) & (0,0228) & (0,0236) & (0,0198) & (0,0239) & (0,0086) & (0,0224) & (0,0122) \\[6pt]
\multicolumn{10}{@{}l}{\textit{R5: especificação de Sakurai (2009) (erro agrupado por município)}} \\
Ano eleitoral & 0,0432\textsuperscript{***} & 0,0689\textsuperscript{***} & $-$0,0257\textsuperscript{***} & 0,2561\textsuperscript{***} & 0,0604\textsuperscript{***} & 0,1909\textsuperscript{***} & 0,0152\textsuperscript{***} & 0,0423\textsuperscript{***} & 0,0227 \\
 & (0,0016) & (0,0020) & (0,0019) & (0,0054) & (0,0025) & (0,0089) & (0,0024) & (0,0059) & (0,0144) \\[6pt]
\multicolumn{10}{@{}l}{\textit{R12: erro-padrão agrupado por município}} \\
Ano eleitoral & 0,0519\textsuperscript{***} & 0,0411\textsuperscript{***} & 0,0343\textsuperscript{***} & 0,2389\textsuperscript{***} & 0,0678\textsuperscript{***} & 0,1778\textsuperscript{***} & 0,0097\textsuperscript{***} & 0,0529\textsuperscript{***} & 0,0318\textsuperscript{**} \\
 & (0,0017) & (0,0021) & (0,0020) & (0,0062) & (0,0029) & (0,0101) & (0,0028) & (0,0069) & (0,0161) \\
Ano pré-eleitoral & 0,0209\textsuperscript{***} & $-$0,0083\textsuperscript{***} & 0,0461\textsuperscript{***} & 0,0509\textsuperscript{***} & 0,0355\textsuperscript{***} & 0,0188\textsuperscript{***} & 0,0077\textsuperscript{***} & 0,0564\textsuperscript{***} & 0,0605\textsuperscript{***} \\
 & (0,0008) & (0,0015) & (0,0011) & (0,0045) & (0,0018) & (0,0072) & (0,0017) & (0,0053) & (0,0122) \\[6pt]
\multicolumn{10}{@{}l}{\textit{R14: sem a dummy de pandemia (erro de Driscoll-Kraay)}} \\
Ano eleitoral & 0,0504\textsuperscript{***} & 0,0649\textsuperscript{***} & $-$0,0078 & 0,2718\textsuperscript{***} & 0,0723\textsuperscript{***} & 0,1964\textsuperscript{***} & 0,0174\textsuperscript{**} & 0,0611\textsuperscript{***} & 0,0434\textsuperscript{***} \\
 & (0,0095) & (0,0198) & (0,0370) & (0,0293) & (0,0119) & (0,0244) & (0,0089) & (0,0168) & (0,0140) \\
Ano pré-eleitoral & 0,0211 & $-$0,0118 & 0,0523\textsuperscript{**} & 0,0461\textsuperscript{*} & 0,0349\textsuperscript{*} & 0,0161 & 0,0065 & 0,0552\textsuperscript{**} & 0,0590\textsuperscript{***} \\
 & (0,0142) & (0,0091) & (0,0249) & (0,0248) & (0,0198) & (0,0233) & (0,0083) & (0,0221) & (0,0122) \\
\hline
\end{tabular}
\fonte{Elaboração própria.}
\nota{Coeficientes do modelo A (efeito fixo de município). Variável dependente e receitas em logaritmo de (1 + valor per capita, R\$ de 2025); o coeficiente de uma dummy se lê como variação percentual aproximada (exata: 100 $\times$ (exp(b) $-$ 1)). Modelo principal e R14: erros-padrão de Driscoll-Kraay; R5 e R12: erros-padrão agrupados por município. O R5 segue \textcite{sakurai2009}: sem ano pré-eleitoral e sem dummy de pandemia; o R14 retira só a dummy de pandemia, e 2020 volta a contar como ano eleitoral. Erros-padrão entre parênteses. \textsuperscript{***}p$<$0,01; \textsuperscript{**}p$<$0,05; \textsuperscript{*}p$<$0,1.}
\end{sidewaystable}

